% ---- pg-040 （书页 28） ----

方法. 此外，还介绍了一个一般性的距离概念，在后面的章节将会反复用到它.

{\large\bfseries 练习}

\mvbold{1.1} 考虑\hyperref[fig:1.1]{图 1.1} 中所画的 7 对测量值 $(x_1, x_2)$：

\begin{mvtable}
\begin{tabular}{lccccccc}
\toprule
$x_1$ & 3 & 4 & 2 & 6 & 8 & 2 & 5 \\
$x_2$ & 5 & 5.5 & 4 & 7 & 10 & 5 & 7.5 \\
\bottomrule
\end{tabular}
\end{mvtable}

计算样本均值 $\bar{x}_1$ 和 $\bar{x}_2$，样本方差 $s_{11}$ 和 $s_{22}$ 及样本协方差 $s_{12}$.

\mvbold{1.2} 一家晨报列出一种外国小轿车的二手车价格，车龄 $x_1$ 以年度量，销售价格 $x_2$ 以千美元度量：

\begin{mvtable}
\begin{tabular}{lccccccccccc}
\toprule
$x_1$ & 1 & 2 & 3 & 3 & 4 & 5 & 6 & 8 & 9 & 11 \\
$x_2$ & 18.95 & 19.00 & 17.95 & 15.54 & 14.00 & 12.95 & 8.94 & 7.49 & 6.00 & 3.99 \\
\bottomrule
\end{tabular}
\end{mvtable}

(a) 作数据的散布图及边缘点图.

(b) 根据散布图推断样本协方差的正负号.

(c) 计算样本均值 $\bar{x}_1$ 和 $\bar{x}_2$ 及样本方差 $s_{11}$ 和 $s_{22}$，计算样本协方差 $s_{12}$ 和样本相关系数 $r_{12}$. 解释这些参数.

(d) 利用\hyperref[eq:1-8]{式(1-8)} 写出样本均值阵列 $\bar{\bm{x}}$，样本方差-协方差阵列 $\bm{S}_n$ 和样本相关阵列 $\bm{R}$.

\mvbold{1.3} 下面是变量 $x_1, x_2$ 和 $x_3$ 的 5 个测量值：

\begin{mvtable}
\begin{tabular}{lccccc}
\toprule
$x_1$ & 9 & 2 & 6 & 5 & 8 \\
$x_2$ & 12 & 8 & 6 & 4 & 10 \\
$x_3$ & 3 & 4 & 0 & 2 & 1 \\
\bottomrule
\end{tabular}
\end{mvtable}

求阵列 $\bar{\bm{x}}, \bm{S}_n$ 和 $\bm{R}$.

\mvbold{1.4} 全球 10 大公司的收益数据如下：

\begin{mvtable}
\mvtabcap{全球 10 大公司 \hfill （10 亿美元）}
\begin{tabular}{lrrr}
\toprule
公司 & $x_1 = $ 销售总额 & $x_2 = $ 利润 & $x_3 = $ 资产 \\
\midrule
花旗集团 & 108.28 & 17.05 & 1 484.10 \\
通用电气 & 152.36 & 16.59 & 750.33 \\
美国国际集团 & 95.04 & 10.91 & 766.42 \\
美洲银行 & 65.45 & 14.14 & 1 110.46 \\
汇丰集团 & 62.97 & 9.52 & 1 031.29 \\
埃克森美孚 & 263.99 & 25.33 & 195.26 \\
皇家荷兰壳牌 & 265.19 & 18.54 & 193.83 \\
英国石油 & 285.06 & 15.73 & 191.11 \\
荷兰国际 & 92.01 & 8.10 & 1 175.16 \\
丰田摩托 & 165.68 & 11.13 & 211.15 \\
\bottomrule
\end{tabular}
\end{mvtable}

{\zihao{6} 资料来源：www. Forbes. com 及福布斯集团的《福布斯》杂志 (2005 年 4 月 18 日).}

(a) 为变量 $x_1$ 和 $x_2$ 作散布图及边缘点图，并加以解释.

(b) 计算 $\bar{x}_1, \bar{x}_2, s_{11}, s_{22}, s_{12}$ 和 $r_{12}$. 解释 $r_{12}$.
